\begin{description}
\item[(1)] Using the \emph{virial theorem} for an isolated, spherical system,
  i.e, that $-2\langle K\rangle = \langle U\rangle$, where ``$K$'' is the
  average kinetic energy and ``$U$'' is the average potential energy of the
  system, determine an expression for the total mass of a cluster of galaxies
  if we know the radial velocity dispersion, $\sigma$, of the cluster's galaxy
  members and the cluster's radius, $R$. Assume that the cluster is isolated,
  spherical, has a homogeneous density and that it consists of galaxies of
  equal mass.
\item[(2)] Find the \emph{virial mass}, i.e.\ the mass calculated from the
  \emph{virial theorem}, of the Coma cluster, which lies at a distance of
  90\,Mpc from us, if you know that the radial velocity dispersion of its
  member galaxies is $\sigma_{v_r} = 1000\,\mathrm{km/s}$ and that its angular
  diameter (on the sky) is about $4^\circ$.
\item[(3)] From observations, the total luminosity of the galaxies comprising
  the cluster is approximately $L = 5 \times 10^{12}\,L_\odot$. If the mass to
  luminosity ratio, $M/L$, of the cluster is ${\sim}1$ (assume that all the
  mass of the cluster is visible mass), this should correspond to a total mass
  $M \sim 5 \times 10^{12}\,\mathrm{M}_\odot$ for the mass of the cluster.
  Give the ratio of the luminous mass to the total mass of the cluster you
  derived in question (2).
\end{description}
